% Abbreviations used in the ECASP thesis.

\abbreviationsAndNotation{
\section*{Abbreviations}
\begin{itemize}
    \item[$-$] \textbf{6DoF:} Six degrees of freedom (three for translation and three for rotation).
    \item[$-$] \textbf{AdaLN-Zero:} Adaptive layer normalization with zero initialization.
    \item[$-$] \textbf{ADD:} Average distance of model points under the estimated and ground-truth poses.
    \item[$-$] \textbf{ADD-S:} Symmetry-aware variant of ADD using closest-point distances.
    \item[$-$] \textbf{BCE:} Binary cross-entropy.
    \item[$-$] \textbf{BOP:} Benchmark for 6D Object Pose Estimation.
    \item[$-$] \textbf{CAD:} Computer-aided design.
    \item[$-$] \textbf{DINO:} Self-distillation with no labels. DINOv3 is used for backbone pretraining.
    \item[$-$] \textbf{DLT:} Direct linear transform, used for triangulation.
    \item[$-$] \textbf{ECASP:} Epipolar Cross-Attention for Silhouette-Based Pose Estimation.
    \item[$-$] \textbf{GELU:} Gaussian error linear unit.
    \item[$-$] \textbf{LoRA:} Low-rank adaptation.
    \item[$-$] \textbf{MLP:} Multilayer perceptron.
    \item[$-$] \textbf{MVPSP:} The calibrated multi-view surgical-instrument pose-estimation benchmark used in this thesis.
    \item[$-$] \textbf{PnP:} Perspective-$n$-Point.
    \item[$-$] \textbf{RGB:} Red, green, and blue color channels.
    \item[$-$] \textbf{RGB-D:} RGB images with depth measurements.
    \item[$-$] \textbf{SfM:} Structure from motion.
    \item[$-$] \textbf{ViT:} Vision Transformer.
    \item[$-$] \textbf{YOLO:} You Only Look Once. YOLO26 is used for instance segmentation.
\end{itemize}

\clearpage
\section*{Mathematical notation}
\begingroup
\renewcommand{\arraystretch}{1.15}
\begin{longtable}{@{}p{0.25\textwidth}p{0.70\textwidth}@{}}
\textbf{Symbol} & \textbf{Meaning} \\
\toprule
\endfirsthead
\multicolumn{2}{@{}l}{\textbf{Mathematical notation (continued)}} \\
\textbf{Symbol} & \textbf{Meaning} \\
\toprule
\endhead
\bottomrule
\endfoot
\(V, v_i, N\) & Set of input camera views, a member view, and the number of views. \\
\(v_0\) & Reference camera view for the object pose and common geometric coordinates. \\
\(R,t\) & Object rotation and translation, mapping object coordinates to the reference camera frame. \\
\(SO(3), SE(3)\) & Three-dimensional rotation group and rigid-transformation group. \\
\(X_{\mathrm{obj}},X_{v_0}\) & A point expressed in the object and reference-camera frames. \\
\(K_v,\tilde K_v\) & Original and crop-adjusted (virtual) camera intrinsic matrices. \\
\(T_{v_i\rightarrow v_j}\) & Rigid transform from camera frame \(v_i\) to camera frame \(v_j\). \\
\(M_v\) & Binary input mask for view \(v\). \\
\(M_v^{\mathrm{vis}},M_v^{\mathrm{full}}\) & Visible and full amodal object masks, respectively. \\
\(\widehat M_v(R,t)\) & Object silhouette rendered in view \(v\) at candidate pose \((R,t)\). \\
\(H_v,W_v\) & Height and width of the original mask in view \(v\). \\
\(B_v,b_w,b_h\) & Foreground crop bounding box and its width and height. \\
\(G,s_v,\delta_u,\delta_v\) & Square canvas size, resize factor, and horizontal/vertical padding offsets. \\
\(A_v,c_v\) & Pixel-coordinate preprocessing transform and normalized crop descriptor. \\
\(\rho,H',W',C\) & Backbone stride, feature-map height and width, and fusion-channel dimension. \\
\(F_v,\widehat F_v\) & Per-view features before and after epipolar cross-view attention. \\
\(o,d,m,\ell\) & Camera center, unit ray direction, Pl\"ucker moment \(m=o\times d\), and concatenated ray code. \\
\(S\) & Number of sampled locations along each epipolar line. \\
\(w_{v,h,w}\) & Normalized aggregation weight of view \(v\) at feature-map location \((h,w)\). \\
\(\pi\) & Perspective projection from a three-dimensional camera point to normalized image coordinates. \\
\(\mathcal L,\lambda\) & Training objective and weighting coefficients of its component losses. \\
\(\mathcal J(R,t)\) & Multi-view silhouette objective minimized during pose refinement. \\
\(\operatorname{cov}_v,\operatorname{exc}_v\) & Silhouette coverage and excess-rendering fractions in view \(v\). \\
\(\varepsilon\) & Positive denominator-stabilization constant. \\
\(\Delta R,\Delta t\) & Rotation error in degrees and translation error in millimeters. \\
\end{longtable}
\endgroup
}
